\cotefile{156-3}
\scanpage{2}
\mspage{1}

\uline{Réseaux}\unsure{\ldots} via découpages.

(8 Juin 86)

% slanted
\margin{\unsure{CF|CH|Cf|Ch}\,\uline{III}}

% slanted
\margin{\deleted{\unsure{Les ensembles}} \unsure{Parties|Ensembles} \deleted{\unsure{des}} de lieux \unsure{modérés|modérées}}

Je voudrais définir une autre axiomatique
\deleted{des réseaux,} \deleted{\unsure{l}}la structure de réseaux
sur un ens. $\unsure{\mathcal L|L}$ (\deleted{\unsure{d}}ens. des ``lieux'').

\textcircled{\scriptsize a}\quad $S_0\deleted{\unsure{(L)|L}}\subset\unsure{P|\mathcal P}_2(L)$ \unsure{l'ens} des \uline{0-sphères}
\uline{plongées}, ou des \uline{paires} \deleted{\unsure{mod.}} (de
lieux) modérées. \inserted{si $\{a,b\}\in S_0$ ; on dit aussi que $a,b$ \unsure{st} \uline{disjoints} ou en position relative \unsure{modérée}.}

% slanted
\margin{Mais, \unsure{on} dit aussi que $a,b$ en position rel. modérée si $a=b$ \uline{ou} $\{a,b\}\in S_0$}

[En termes de
l'axiomatique par les segments,
cel\unsure{a|à|} signifie que \deleted{$a,b$ \unsure{sont} ou bien} \inserted{(à \unsure{part} $a,b$ \uline{distincts}},
\unsure{contenu} dans $\{a,b\}\in\deleted{\unsure{S_0}}P_2(\mathcal L)$, que
ou bien $\{a\unsure{;|,}b\}\cap\mathcal L_s\ne\varnothing$ ($a$ ou $b$ est un
nœud \unsure{sg$_r$|sg} du réseau) \uline{ou} \uline{bien} \deleted{$a,b$}
(si \deleted{\unsure{st|sont}} \inserted{$a,b$} réguliers) que $a,b$ \inserted{\unsure{doivt|doit}} \unsure{appartiennent|appartient} :
des lignes différentes, \unsure{soit|sont|sauf}, s'ils sont
dans une $\unsure{\bar{=}|=}$ ligne, \unsure{pour|que} $\{a,b\}=\partial I$,
pour $I\in S$].

\deleted{\unsure{b)}} Une partie \uline{finie} \inserted{$T$} est dite \deleted{\unsure{admissible}} \inserted{``\uline{modérée}''}
si \deleted{$\forall\,a,b\in T$, \unsure{on a} \unsure{$a,b$}} $P_2(T)\subset S_0$
i.e. deux \unsure{éléments|élémts|elements} quelc. de $T$ sont en
position relative modérée (\deleted{\unsure{des}} \unsure{2 st|sont} ``disjoints'').

% paper watermark
\rotatebox{90}{\textsc{\unsure{x}tra Strong Ext}}
\scanpage{3}
\mspage{2}

\deleted{De plus, on veut que la structure \unsure{implique}}

\deleted{\textcircled{\scriptsize b}\quad $\forall\,T\in\unsure{P|\mathcal P}_{f\,\mathrm{mod}}(\unsure{\mathcal L|L})$, on se donne une
relation d'équivalence $\unsure{\mathcal R|R}_T$ dans l'ens
\unsure{où}}

Si $T$ est une partie finie
modérée\unsure{,|;}, on \deleted{\unsure{désignera|désigne} par} \inserted{pose}

\[
 C\deleted{\unsure{\mathrm{omp}^{r}|\mathrm{omp}}}_{\mathcal L}(T)=
 \{\,t\in\mathcal L\setminus T\mid \{T\cup\{t\}\}\ \text{\unsure{modéré|modérée}}\,\}
\]

% slanted
\margin{\uline{le complémentaire} \uline{modéré} de $T$.
Question : $T$ est-il connu quand on \unsure{connait|connaît} $C_{\mathcal L}(T)$ ? Plus précisément,
si $T'\subset T$, on a $C(T)\subset C(T')$, l'inverse est-il vrai ? Oui,
Donc, si \inserted{$t\in T'$} $t\notin T$, \deleted{\unsure{$t$ est gén.} \unsure{tel} ($T\cup\{t\}$ modéré) \unsure{puisque} \unsure{on a toujours} \unsure{$s\notin T$,}} alors $t\in C(T)$ et $t\notin C(T')$}

\unsure{ou mieux|ou même}, l'ens des $t\in L\setminus T$ en
position modérée p.r. à $T$.

[\uline{NB}\quad deux parties finies modérées
sont dites en position relative
modérée si leur\unsure{-| }réunion est
modérée ; $t$ est dit en pos. rel.
modérée p.r. à $T$ si $\{t\}$ l'est ;
$s$ et $t$ sont dits en pos. rel. modérée
si $\{s,t\}$ est modéré i.e. \deleted{\unsure{si}}\unsure{,}
$s\ne t$ ou $\{s;t\}\in S_0$ \ldots\ ]

% slanted
\margin{\uline{NB} $T$\deleted{\unsure{est}} est connu quand on \unsure{connait|connaît} $C(T)$}

\textcircled{\scriptsize b}\quad $\forall\,T\in\underbrace{\mathrm{Mod}_{f}}_{\substack{\text{parties mod\unsure{.|}}\\\text{finies}}}(\mathcal L)$, une relation
d'équivalence $\mathcal R_T$ dans $C_{\unsure{\checkmark|\mathcal L|}}(T\unsure{)|\}}$ (dont
les classes s'appellent les \uline{comp. connexes} \inserted{de $C(T)$}

% paper watermark, visible from the reverse
\rotatebox{90}{\textsc{\unsure{x}tra Strong Ext}}
\scanpage{4}
\mspage{3}

\textcircled{\scriptsize c}\quad $\forall\,T\in\mathrm{Mod}_{f}(\unsure{\mathcal L|L})$, et $A\in C(T)/\unsure{R|\mathcal R}_T$
une classe suivant $R_T$, une partie
\unsure{de} $\dot A\subset T$, appelée \uline{frontière} de $A$.

\unsure{Ces|Les} données \textcircled{\scriptsize a} \textcircled{\scriptsize b}, \textcircled{\scriptsize c} étant \unsure{assujetties}
à satisfaire certains axiomes.

Je prévois que les axiomes impliquent
\deleted{Soit \illegible{1} $\unsure{t|U}\in C(T)$, \unsure{on pose}}
ceci

1°) que la donnée \textcircled{\scriptsize c} résulte déjà de \textcircled{\scriptsize b}

2°) que la donnée \textcircled{\scriptsize b} résulte de
la donnée des $R_T$, \deleted{\unsure{pour}} dans le cas
particulier où $T\in S_0$. Donc
les seules \unsure{donnée|données} essentielles \unsure{serait|seraient}
\textcircled{\scriptsize a} \unsure{celles|celle} des 0-sphères i.e. la notion
de \deleted{\unsure{pos.}} lieux en pos. relative modérée
(\unsure{relation|condition} \inserted{réflexive et} symétrique) et \textcircled{\scriptsize b} La donnée,
pour \deleted{\unsure{tout}} $\unsure{\varepsilon|\epsilon|E}\in S_0$, de la relation
d'équivalence $R_{\varepsilon}$ \unsure{dans} $C(\varepsilon)$.

\deleted{\unsure{On va|Donc on}} \uline{Axiomes provisoires :}

\deleted{Soit \unsure{$X$}} \unsure{Appellons|Appelons} les $U$ de la forme
$C(\deleted{\unsure{L}}T)$ \uline{parties modérées cofinies}. On

% paper watermark
\rotatebox{90}{\textsc{\unsure{S}trong Extra St}}
\scanpage{5}
\mspage{4}

veut que si $U$ \deleted{$(=C\unsure{)(|(}T)$} est à la fois modéré fini
et modéré cofini, donc si $S$ et $T$
sont \deleted{\unsure{deux}} parties modérées finies
\deleted{\unsure{telles}} disjointes telles que \deleted{\unsure{tout}} \inserted{a)} $S\cup T$
modéré i.e. \deleted{\unsure{tout}} $\unsure{s|S}\in S$ modéré p.r. à \unsure{tt}
$t\in T$, et b) Tout $u\in\unsure{\mathcal L|L}$ \unsure{modér.|modéré} p.r. à
$S$ est dans $T$\unsure{;|,}, \deleted{alors $S\cup T=\mathcal L$} et
à fortiori $S\cup T$ est partie \inserted{finie} modérée
maximale, alors $S\cup T=\mathcal L$, et
$\unsure{{}^{\cdot\cdot}S_0|S_0}=\unsure{P|\mathcal P}_2(\mathcal L)$, enfin les relations
d'équivalence sont discrètes, et les
$\unsure{\dot A|\bar A}$ sont vides. Pour \unsure{le|la} formuler,
on définit \unsure{la|a}

\deleted{\unsure{Déf}} \inserted{Ex 1} \uline{structure \inserted{de \unsure{Réseau|réseau}} discrète} sur $\mathcal L$

\[
 S_0=P_2(\mathcal L)\ \Longrightarrow\
 \left\{\begin{aligned}
 &\mathrm{Mod}_f(\unsure{\mathcal L|\mathcal L_0})=\deleted{\unsure{\mathrm{Mod}}}P_f(\mathcal L)\\
 &C_T=\mathcal L\setminus T\qquad\text{si }T\in P_f(\mathcal L)
 \end{aligned}\right.
\]

\[
\begin{aligned}
 &\deleted{C(T)=\unsure{\subset|\mathcal L}}\\
 &\unsure{R|\mathcal R}_{\unsure{T|T\rceil|\Pi}}=\text{relation d'équivalence \unsure{discrète|discrèt}}\\
 &\hspace{3em}\text{dans }\mathcal L\setminus T,\quad\text{donc les \unsure{composantes}}\\
 &\hspace{4em} \inserted{A}\ \text{sont réduites à un pt}\\
 &\deleted{\text{\unsure{Pour Donc}}}\\
 &\dot A=\varnothing
\end{aligned}
\]

\uline{Axiome}\quad Si $T$ est partie finie modérée
maximale, alors $T=\mathcal L$, $\mathcal L$ est \unsure{discrèt|discrète|discret}.

% paper watermark, visible from the reverse
\rotatebox{90}{\textsc{\unsure{S}trong Extra St}}
\scanpage{6}
\mspage{5}

On va faire des \unsure{hyp.thèses|hypothèses} de
finitude \unsure{draconiennes|draconniennes}, peut-être
provisoires

% slanted
\margin{Finitude}

$\underline{\mathrm{Dec}}_1$\quad \deleted{\unsure{l}}l'ens des comp. connexes
de tout $U$ est \uline{fini}.

\deleted{\uline{Définition} Partie modérée \uline{fermée}}

$\underline{\mathrm{Dec}}_2$\quad \unsure{Soit|Soient}\deleted{\unsure{ent}} $U$\inserted{$=C(\deleted{\unsure{T}}S), V=C(\deleted{\unsure{S}}T)$} \unsure{partie|parties} modérée
cofinies, avec $U\subset V$ i.e. $T\subset S$.
Alors \unsure{tte} composante connexe \inserted{$A$} de
$U$ est contenue dans une comp.
connexe (unique bien \unsure{sur|sûr}) $B$ de $V$.

De plus, on a

\diagram{\begin{tikzcd}[ampersand replacement=\&,row sep=0.35ex,column sep=0.4em,
cells={nodes={inner sep=1pt}},execute at end picture={
\draw[line width=.35pt] ([yshift=-.2ex]aa.south west) to[out=-90,in=180] ([xshift=.3ex,yshift=-.55ex]aa.south west) -- ([xshift=-.3ex,yshift=-.55ex]aa.south east) to[out=0,in=-90] ([yshift=-.2ex]aa.south east);
\draw[line width=.35pt] ([yshift=-.2ex]bb.south west) to[out=-90,in=180] ([xshift=.3ex,yshift=-.55ex]bb.south west) -- ([xshift=-.3ex,yshift=-.55ex]bb.south east) to[out=0,in=-90] ([yshift=-.2ex]bb.south east);
}]
|[alias=aa]| \dot A \& \cap\,T\,\subset \& |[alias=bb]| \dot B\\
\cap \& \& \cap\\
\unsure{S|s} \& \& T
\end{tikzcd}}

% paper watermark; read from foot to head; clipped at the two sheet edges
\rotatebox{90}{\textsc{\unsure{t}rong Extra Str}}
\scanpage{7}
\mspage{6}

\uline{Ex \deleted{\unsure{1}}\unsure{2|1}}\quad Soit $\deleted{\unsure{S|E}}$ $\unsure{\mathcal L|L}$ un ens ordonné,
on suppose $\mathcal L$ filtrant croissant, filtrant
décroissant, sans plus grand ni
plus petit \unsure{élémt|élément}, loc. filtrant croissant
et filtrant décroissant, divisible.

On \unsure{appelera|appellera} un tel ensemble une
\uline{ligne ordonnée}

\[
S_0\subset\unsure{P|\mathcal P}_2(\mathcal L)\qquad
\deleted{\text{\unsure{Soit} }\{x,y\}\in P_2(\mathcal L)}
\]

\textquotedbl\quad \unsure{partie|parties} \deleted{qui} de card. 2 \deleted{\unsure{une}} qui sont \uline{tot.} ordonnées

\[
\begin{aligned}
&\mathrm{Mod}_f(\mathcal L)=\text{parties finies \unsure{tot|tot.} ordonnées}\\
&\hspace{8em}\deleted{\text{i.e. telles que}}
\end{aligned}
\]

\[
\begin{aligned}
 \text{Soit }T\in\mathrm{Mod}_f(\mathcal L),\qquad &T=\{t_{\unsure{p|1}},t_2,\ldots,t_{\unsure{n|m}}\}\\
 &\quad t_1<\unsure{t|l|1}_2<\ldots<t_n
\end{aligned}
\]

Alors les composantes connexes sont

\[
\underbrace{\mathopen{]}-\infty,t_1\mathclose{[}}_{t_1},\
\underbrace{\mathopen{]}t_1,t_2\mathclose{[}}_{t_1,t_2},\
\ldots\quad
\underbrace{\mathopen{]}t_{n-1},t_n\mathclose{[}}_{t_{n-1}\,t_n},\
\underbrace{\mathopen{]}t_n,\infty\mathclose{\unsure{]|[}}}_{t_n}
\]

% slanted
\margin{\unsure{pts|Pts} frontière}

($n\unsure{+|-}1$ \unsure{composantes} connexes)

% paper watermark; mirror impression across the letters
\rotatebox{90}{\textsc{\unsure{t}rong Extra Str}}
\scanpage{8}
\mspage{7}

Cette structure ne change pas si
\deleted{\unsure{de}}on change l'ordre \unsure{par} ordre opposé,
i.e. elle ne dépend que du
biordre.

\deleted{Partie modérée} \deleted{\unsure{Segments|Segment}}

% slanted
\margin{refaire, sans se borner à une ligne : \unsure{sous tronçon|son tronçon} fermé d'un tronçon. C'est un tronçon \unsure{p.|pour} le biordre induit. D'autre part, le bord d'un tronçon est défini \unsure{comme|Ens} des}

\unsure{Tronçons|Tronçon} de la ligne : parties
de la forme

\[
\left\{\begin{aligned}
 &[s,t] &&s\leq t\\
 &\mathopen{]}-\infty,s],\quad\text{ou}\quad [s,+\infty\mathclose{[}\\
 &\unsure{\mathcal L|L}=\mathopen{]}-\infty,+\infty\mathclose{[}
\end{aligned}\right.
\]

Bord du tronçon : $\{s,t\}$ dans le
premier cas, $\{s\}$ dans le second, $\varnothing$ \inserted{troisième}

Deux tronçons sont dits en
position relative modérée, si leurs
bords \unsure{:|} \unsure{le|les} sont.

\uline{Partie} \uline{modérée} :
réunion \inserted{finie} de tronçons \inserted{en position} mutuelle
modérée, et disjoints. \deleted{\unsure{Donc}} Par
\unsure{définition|definition}, son bord est la réunion
(disjointe) des bords.

% paper watermark; clipped at the two sheet edges
\rotatebox{90}{\textsc{\unsure{r}ong Extra Stro}}
\scanpage{9}
\mspage{8}

\deleted{\uline{Proposition}}

\deleted{\uline{Définition} $\overbrace{\text{Segment, intervalle, demi-intervalle}}^{\inserted{\text{tronçons}}}$
(comme structure de réseau).}
\deleted{Segment}

\uline{Définition}\quad Un ens. ordonné \inserted{$\unsure{\mathcal L|L}$} est appelé un \uline{tronçon} \uline{ordonné} si

\unsure{a}) $\forall\,a\in\mathcal L$ tel que $\mathcal L_{>a}\deleted{\unsure{\sim|=}}\unsure{\ne|=}\varnothing$, $\mathcal L_{>a}$ est filtrant décroissant
\deleted{i.e. $\forall\,x,y$ $\unsure{a|x},y$ tels que} \inserted{si $a,y,a\in\mathcal L$ tels que}
$a>a$, $y>a$ alors $\exists\,z\in\mathcal L$ avec $x,y\geq z>a$

$\deleted{\unsure{b|a}}a'$) Axiome dual : si $a,x,y\in\mathcal L$ tels \unsure{q|que}
$\unsure{x|a}<a$, $y<a$ alors $\exists\,z\in\mathcal L$ avec $a,y\leq z<a$

b) $\mathcal L$ est filt\unsure{ant|rant} croissant et filtrant décroissant i.e.
$x,y\in\mathcal L\Rightarrow\exists\,a,b\in\mathcal L$ avec $a\leq x,y\leq b$

c) Si $x,y\in\mathcal L$ tels que $x<y$, $\exists\,z$ tel que $x<z<y$ .

d) $\deleted{\mathcal L\,\unsure{\ne|=}}$ $\operatorname{Card}\mathcal L>1$ \deleted{\unsure{(imp.|(}}

% slanted
\margin{\uline{NB} les cond\unsure{.|} $a\ a'\ b)\ c)$ sont satisfaites si $\operatorname{card}\mathcal L\leq1$ !}

On dit que $\mathcal L$ est un \uline{segment} si de plus $\mathcal L$ admet un plus petit et
\deleted{\unsure{(p.|(}} un plus grand élément, un \uline{intervalle} s'il n'admet ni l'un ni l'autre,

% paper watermark, visible from the reverse
\rotatebox{90}{\textsc{\unsure{n}g Extra Stro}}
